% !TeX program = xelatex
\documentclass[UTF8]{ctexart}

% -----------------------------------------------------------
% 1. 引入宏包
% -----------------------------------------------------------
\usepackage[a4paper, margin=2cm]{geometry}
\usepackage{tcolorbox}
\usepackage{amsmath, amssymb, amsfonts}
\usepackage{booktabs}
\usepackage{xcolor}
% 引入 skins 库以支持高级皮肤和绘图
\tcbuselibrary{listings, skins, breakable}

% -----------------------------------------------------------
% 【全局排版优化】
% -----------------------------------------------------------
% 教程类文档通常不需要首行缩进，而是靠段落间距区分
\setlength{\parindent}{0pt}  % 全局取消首行缩进
\setlength{\parskip}{8pt}    % 增加段落之间的垂直间距

% -----------------------------------------------------------
% 2. 定义深色主题颜色
% -----------------------------------------------------------
\definecolor{editorDark}{RGB}{40, 44, 52}     % 代码背景深灰色
\definecolor{editorFrame}{RGB}{33, 37, 43}    % 边框更深
\definecolor{editorTitle}{RGB}{171, 178, 191} % 标题文字灰白色
\definecolor{keywordBlue}{RGB}{97, 175, 239}  % 关键字亮蓝
\definecolor{cmdPurple}{RGB}{198, 120, 221}   % 命令亮紫
\definecolor{strGreen}{RGB}{152, 195, 121}    % 字符串亮绿

% -----------------------------------------------------------
% 3. 配置深色模式代码高亮
% -----------------------------------------------------------
\lstdefinestyle{darkstyle}{
	language=TeX,
	basicstyle=\ttfamily\small\color{white},
	keywordstyle=\color{keywordBlue}\bfseries,
	commentstyle=\color{gray!60}\itshape,
	texcsstyle=*\color{cmdPurple},
	morekeywords={begin, end, section, subsection, textbf, textit, underline, mathbb, mathcal},
	numbers=left,
	numberstyle=\tiny\color{gray!50},
	frame=none,
	breaklines=true,
	columns=flexible,
	keepspaces=true,
	tabsize=4,
	showtabs=false
}

% -----------------------------------------------------------
% 4. 设计 macOS 窗口风格 (修复了空行导致的编译错误)
% -----------------------------------------------------------
\tcbset{
	macos-editor/.style={
		skin=bicolor,
		enhanced, % 显式启用增强模式，确保 title 节点可用
		colframe=editorFrame,
		colback=editorDark,
		colbacklower=white,
		coltitle=editorTitle,
		colupper=white,
		arc=6pt,
		boxrule=1pt,
		drop fuzzy shadow,
		left=4pt, right=4pt, top=2pt, bottom=2pt,
		fonttitle=\sffamily\bfseries\small,
		toptitle=5pt, bottomtitle=5pt,
		lefttitle=45pt, % 左侧留白给按钮
		overlay={
			\begin{scope}[shift={(title.west)}]
				\fill[red!60]   (15pt, 0) circle (3pt);
				\fill[orange!60] (28pt, 0) circle (3pt);
				\fill[green!60]  (41pt, 0) circle (3pt);
			\end{scope}
		}
	}
}

% -----------------------------------------------------------
% 5. 定义环境：左代码，右结果
% -----------------------------------------------------------
\newtcblisting{code-left}[2][]{
	macos-editor,
	listing style=darkstyle,
	sidebyside,
	listing and text,
	lefthand width=0.55\textwidth,
	valign=center,
	title={#2},
	#1
}

% -----------------------------------------------------------
% 6. 定义环境：上代码，下结果
% -----------------------------------------------------------
\newtcblisting{code-top}[2][]{
	macos-editor,
	listing style=darkstyle,
	title={#2},
	#1
}

% -----------------------------------------------------------
% 正文开始
% -----------------------------------------------------------
\begin{document}
	
	\begin{center}
		{\Huge \bfseries LaTeX 极简教程} \\
		\vspace{0.3cm}
		{\large Dark Mode Edition (macOS Style)}
	\end{center}
	
	\tableofcontents
	\newpage
	
	% ==================================================
	\section{文本格式 (Text Formatting)}
	% ==================================================
	
	\begin{code-left}{字体样式演示}
		% 常用字体命令
		\textbf{Bold (加粗)} \\
		\textit{Italic (斜体)} \\
		\underline{Underline (下划线)} \\
		\texttt{Monospace (等宽)} \\
		\textsc{Small Caps (小型大写)}
	\end{code-left}
	
	\begin{code-left}{字号大小对比}
		{\tiny tiny (极小)} \\
		{\footnotesize footnote (脚注)} \\
		{\small small (小号)} \\
		{\normalsize normal (正常)} \\
		{\Large Large (大号)} \\
		{\huge huge (巨大)}
	\end{code-left}
	
	% ==================================================
	\section{数学公式 (Mathematics)}
	% ==================================================
	
	\subsection{基础公式}
	
	\begin{code-left}{行内与行间公式}
		% 行内公式使用 $...$
		质能方程：$E = mc^2$
		
		% 行间公式使用 \[...\]
		薛定谔方程：
		\[
		i\hbar\frac{\partial}{\partial t}\Psi = \hat{H}\Psi
		\]
	\end{code-left}
	
	\subsection{复杂结构}
	
	% 使用 code-top 展示较长的代码
	\begin{code-top}{分段函数 (Cases)}
		% 需要 amsmath 宏包
		定义符号函数：
		\[
		f(x) = 
		\begin{cases} 
			-1 & \text{if } x < 0 \\
			0  & \text{if } x = 0 \\
			1  & \text{if } x > 0
		\end{cases}
		\]
	\end{code-top}
	
	\begin{code-left}[lefthand width=0.6\textwidth]{矩阵示例}
		% 各种矩阵括号
		Pmatrix: $\begin{pmatrix} a & b \\ c & d \end{pmatrix}$ \\[5pt]
		Bmatrix: $\begin{bmatrix} 1 & 2 \\ 3 & 4 \end{bmatrix}$ \\[5pt]
		Vmatrix: $\begin{vmatrix} x & y \\ z & w \end{vmatrix}$
	\end{code-left}
	
	% ==================================================
	\section{表格与列表}
	% ==================================================
	
	\begin{code-top}{三线表示例 (Booktabs)}
		% 专业的学术表格
		\begin{tabular}{llr}
			\toprule
			Name & Role & Salary \\
			\midrule
			Alice & Engineer & \$80k \\
			Bob   & Manager  & \$100k \\
			\bottomrule
		\end{tabular}
	\end{code-top}
	
\end{document}